\section{Thảo luận}
Kết quả thử nghiệm cho thấy kiến trúc đề xuất phù hợp với bài toán giám sát và điều khiển môi trường trong căn hộ. So với hướng chỉ dùng thiết bị lọc không khí hoặc chỉ giám sát cây trồng, hệ thống kết hợp cả dữ liệu IAQ và dữ liệu chăm sóc cây, từ đó hỗ trợ người dùng nhìn nhận môi trường sống theo một bức tranh đầy đủ hơn. Việc sử dụng LoRa cho kết nối node - gateway giúp node không phụ thuộc trực tiếp vào Wi-Fi, phù hợp với vị trí đặt cây hoặc khu vực sóng Wi-Fi yếu. ThingsBoard giúp rút ngắn thời gian phát triển nhờ có sẵn chức năng lưu trữ telemetry, dashboard và điều khiển từ xa. Flutter phù hợp để xây dựng ứng dụng di động có giao diện trực quan và dễ mở rộng.

Một số hạn chế vẫn cần được cải thiện. Thứ nhất, độ chính xác của cảm biến TDS, pH và CO2 phụ thuộc nhiều vào hiệu chuẩn; đây cũng là thách thức thường gặp trong các hệ thống IoT môi trường. Thứ hai, điều khiển điều hòa bằng IR là điều khiển một chiều; node cần có cơ chế lưu trạng thái dự đoán hoặc bổ sung cảm biến phản hồi để biết chắc điều hòa đã nhận lệnh. Thứ ba, hệ thống mới triển khai với một node nên chưa đánh giá đầy đủ khả năng mở rộng nhiều node. Thứ tư, đề tài chưa đo trực tiếp VOC hoặc bụi mịn, nên chưa thể kết luận định lượng đầy đủ về khả năng lọc không khí của cây. Các hạn chế này là cơ sở cho hướng phát triển tiếp theo.
